\section{Discussion and limitations}
\label{sec:disc}
\paragraph{What the evidence supports.}
On German Credit (four models) and ACS Employment (\numacsEmpModels), a single Kernel SHAP run at the library default is adequate for most instances and its noise is smaller than the change that a retrained model brings under our retraining protocol; this is two-dataset evidence on one trained model per cell; the cases to worry about are instances with a small margin between the last attribute shown and the first omitted, and any setting below the default budget. For $d\le11$ the default is exact; in our stress-test grid below the default, top-$k$ stability degrades while the error still looks small, but this depends on where our grid starts. These are statements about the cells we ran (Table~\ref{tab:setup}); they are not a guarantee for other datasets, models or $d$.
\paragraph{What it means for practice.}
If an explanation is quoted as a top-$k$ set, report the margin or repeat a small number of runs on the instance; for Kernel SHAP, averaging $m$ runs was not a better use of compute than one run $m$ times larger, and for permutation SHAP the two were indistinguishable (Section~\ref{sec:exp}). Check the library defaults: in the cells where we tested them (a subset), \texttt{l1\_reg} and the explained output changed the result by more than the seed did.

\subsection{Limitations}
\label{sec:limits}
\begin{itemize}
\item \textbf{Scope of inference.} One trained model per cell and 50 explained instances per cell: statements are conditional on that model and apply to instances of it. Five datasets (ACS: California 2018 only; two are tiny), four model classes; no deep or very high-dimensional data.
\item \textbf{Regime.} For $d\le10$ we study sub-default budgets because the library default is exact; the default-budget claims rest on the cells with $d>10$ only.
\item \textbf{Repeat versus enlarge.} The comparison exists only at the budget pairs of our grid whose compute ratio is within $[0.9,1.1]$ for $c=2,4,8$ repeats (Section~\ref{sec:exp}), has no matched permutation pairs at $c=8$, and its intervals rest on five datasets; we make no statement for other ratios.
\item \textbf{Dependence between cells.} Seeds are the same function of (instance, budget index, run) in every cell, so the cells of one dataset sample the same coalitions for the same instances; the rank correlation of run noise between two models of one dataset is on average \numseedCorrMeanK{} (Kernel) and \numseedCorrMeanP{} (permutation), at most \numseedCorrMaxK{}, against about zero with shuffled run index. The effective number of independent units is closer to the five datasets than to the \numnCells{} cells.
\item \textbf{Retrain standard.} One protocol (80\% subsamples, five seeds, one split); the subsample fraction was varied (50\%, 95\%) only for 11 cells (appendix), and the retrained models share most of their data. The background standard is reduced (three redraws) for $d>10$.
\item \textbf{Reference.} For German Credit the reference is itself an estimate (the mean of 10 runs at 20\,000 coalitions); it was checked against enumeration on \numrefN{} instances of two models, which supports ``no gross bias'' and nothing more.
\item \textbf{Library-specific.} All results are for \texttt{shap} 0.51.0 with its paired sampling and its seeding behaviour; the default \texttt{l1\_reg} and log-odds checks cover a subset of cells (Section~\ref{sec:exp}).
\item \textbf{Not done.} The leverage-score estimator of \citet{musco2024leverage}, a tree-based exact comparison, the Covert--Lee variance diagnostic \citep{covert2020improving}, the Arrhythmia stress test and German Credit budgets beyond 5000 were not run, and the per-instance diagnostic of Section~\ref{sec:diag} is exploratory.
\end{itemize}
